\documentclass[../paper.tex]{subfiles}
\begin{document}
\section{Conclusion}\label{sec:conclusion}
This paper develops a model of simultaneous commitment, focusing on the possibility that commitment can lead to worse outcomes. In the model, players can voluntarily commit to normal-form strategies before playing a base game. We show that in some symmetric $2 \times 2$ games, the option to commit strictly reduces equilibrium payoffs. We characterize symmetric $2 \times 2$ games for which commitment reduces or improves equilibrium payoffs and extend the model to multiple commitment periods. As $T \to \infty$, early incompatible commitments can become more likely, and payoffs can decline further. 
 
\paragraph{Several limitations remain.}
First, the finding that commitment can worsen outcomes relies on mixed-strategy Nash equilibria as the equilibrium concept. We focus on mixed-strategy equilibria because they capture the possibility of miscoordination and allow for quantifying its likelihood. Moreover, in symmetric games, mixed equilibria may be the only symmetric equilibria available. However, this justification does not extend to asymmetric games.
 
Secondly, this paper considers only a narrow class of $2\times2$ games. Analyzing a broader class of games may reveal additional dynamics.
 
Finally, while motivated by concerns over AI agents capable of making commitments, this paper does not explore technical implementation. An interesting direction for further research is to explore in what ways commitment in AI agents can be designed to avoid the inefficient outcomes identified in this paper. Conditional commitments \parencite[]{kalai_commitment_2010} or simulation-based program equilibria \parencite[]{cooper_characterising_2024} seem like promising directions here. 
 
\newpage
\printbibliography
\newpage

\end{document}